% !TeX root = ../../elegantbook-cn.tex


\section{条件翻译}
\subsection{坐标}
\noindent（1）直线与$y=\frac{y_1}{x_1 + 2}(x + 2)$或者$x = 4$交于点$P$，求点$P$坐标 \\
（2）直线$y = \frac{y_0}{x_0}x$与直线$y = \frac{2(2y_0 - 3)}{x_0}x + 3$交于点$Q$，求点$Q$坐标 \\
（3）过点$A(-2, 0)$斜率为$k$的直线与椭圆$\frac{x^2}{4} + y^2 = 1$交于另一点$B$，求点$B$坐标 \\
（4）中点坐标公式、与椭圆相关的中点处理\\
举个例子：$\frac{x^2}{4} + y^2 = 1$，直线$y = x+ 1$上有一点$M$，射线$OM$与椭圆交于点$N$，若$M$是$ON$的中点，求点$M$的坐标

\begin{example} % 设线、中点处理
（2024西城期末）已知椭圆 $E:\frac{x^{2}}{a^{2}}+\frac{y^{2}}{b^{2}}=1\ (a>b>0)$ 的离心率为 $\frac{\sqrt{3}}{2}$，且经过点 $C(2,1)$.\\
(1) 求椭圆 $E$ 的方程；\\
(2) 过点 $N(0,1)$ 的直线交椭圆 $E$ 于点 $A,B$（点 $A,B$ 与点 $C$ 不重合）. 设 $AB$ 的中点为 $M$，连接 $CM$ 并延长交 $E$ 于点 $D$. 若 $M$ 恰为 $CD$ 的中点，求直线 $AB$ 的方程.

\end{example}

\begin{example} % 直线，与椭圆相关的中点处理
（2023西城期末）如图，已知椭圆 $E:\frac{y^2}{a^2}+\frac{x^2}{b^2}=1\ (a>b>0)$ 的一个焦点为 $F_1(0,1)$，离心率为 $\frac{\sqrt{2}}{2}$.\\
（1）求椭圆 $E$ 的方程；\\
（2）过点 $F_1$ 作斜率为 $k$ 的直线交椭圆 $E$ 于两点 $A,B$，$AB$ 的中点为 $M$. 设 $O$ 为原点，射线 $OM$ 交椭圆 $E$ 于点 $C$. 当 $\triangle ABC$ 与 $\triangle ABO$ 的面积相等时，求 $k$ 的值.
\end{example}


\subsection{距离公式}
\noindent （1）点到点的距离公式：点$A(x_1, y_1)$与点$B(x_2, y_2)$之间的距离为$\sqrt{(x_2 - x_1)^2 + (y_2 - y_1)^2}$ \\
（2）弦长公式：直线$y = kx + m$上两点$A(x_1, y_1), B(x_2, y_2)$之间的距离为$d = \sqrt{1 + k^2} \cdot \vert x_2 - x_1\vert = \sqrt{1 + \frac{1}{k^2}} \cdot \vert y_2 - y_1\vert$\\；直线$x = ny + t$上两点$A(x_1,y_1)$与点$B(x_2, y_2)$之间的距离为$d = \sqrt{1 + n^2}\vert y_1 -y_2 \vert$。弦长公式常用于直线斜率已知以及求同一条直线上两线段的长度的比值\\
（3）点到直线的距离公式：点$P(x_0, y_0)$到直线$Ax + By + C = 0$的距离为$d = \frac{|Ax_0 + By_0 + C|}{\sqrt{A^2 + B^2}}$

\begin{example} % 动点，直线相交求坐标，距离公式
（2026丰台期末）已知 $F_1,F_2$ 是椭圆 $E:\frac{x^2}{a^2}+\frac{y^2}{b^2}=1\ (a>b>0)$ 的左、右焦点，且 $\lvert F_1F_2\rvert=2$，点 $Q\left(1,\frac{3}{2}\right)$ 在椭圆 $E$ 上.\\
（1）求椭圆 $E$ 的标准方程；\\
（2）设 $P(x_0,y_0)$ 为第一象限内 $E$ 上的动点（不与 $Q$ 重合），直线 $PF_1,PF_2$ 分别交直线 $l:4y_0x-3x_0y=0$ 于点 $M,N$. 求 $\lvert MF_1\rvert+\lvert NF_2\rvert$ 的值.
\end{example}
\begin{example} % 直线或者动点，距离公式，等腰直角三角形
（2026海淀期末）已知椭圆 $E:\frac{x^2}{a^2}+\frac{y^2}{b^2}=1\ (a>b>0)$ 的左顶点为 $A$，请从下面条件中任选两个作为已知.\\
条件①：椭圆 $E$ 的离心率为 $\frac{1}{2}$；\\
条件②：椭圆 $E$ 经过点 $(0,\sqrt{3})$；\\
条件③：点 $A$ 到椭圆 $E$ 的左焦点的距离为 $1$.\\
（1）求椭圆 $E$ 的方程；\\
（2）过点 $A$ 的直线 $l$ 与椭圆 $E$ 的另一个交点为 $P$，$l$ 与 $y$ 轴交于点 $M$. 过点 $P$ 作 $l$ 的垂线交直线 $x=2$ 于点 $N$. 若 $\triangle PMN$ 是等腰直角三角形，求 $l$ 的斜率.
\end{example}

\begin{example} % 动点引入，距离公式
（2024海淀一模）已知椭圆 $G:x^2+my^2=m$ 的离心率为 $\frac{\sqrt{2}}{2}$，$A_1,A_2$ 分别是 $G$ 的左、右顶点，$F$ 是 $G$ 的右焦点.\\
（1）求 $m$ 的值及点 $F$ 的坐标；\\
（2）设 $P$ 是椭圆 $G$ 上异于顶点的动点，点 $Q$ 在直线 $x=2$ 上，且 $PF\perp FQ$，直线 $PQ$ 与 $x$ 轴交于点 $M$. 比较 $\lvert MP\rvert^2$ 与 $\lvert MA_1\rvert\cdot\lvert MA_2\rvert$ 的大小.
\end{example}

\begin{example} % 动点引入（入题思路比较难），垂直的翻译，弦长公式
（2021海淀二模）椭圆 $C:\frac{x^2}{a^2}+\frac{y^2}{b^2}=1\ (a>b>0)$ 的左、右焦点分别为 $F_1,F_2$，$E$ 是椭圆 $C$ 上一点，且 $\lvert F_1F_2\rvert=2$，$\lvert EF_1\rvert+\lvert EF_2\rvert=4$。\\
（1）求椭圆 $C$ 的方程；\\
（2）$M,N$ 是 $y$ 轴上的两个动点（点 $M$ 与点 $E$ 位于 $x$ 轴的两侧），$\angle MF_1N=\angle MEN=90^\circ$，直线 $EM$ 交 $x$ 轴于点 $P$，求 $\frac{\lvert EP\rvert}{\lvert PM\rvert}$ 的值。
\end{example}






\subsection{斜率和角度}
\noindent（1）两直线平行：直线斜率$k_1 = k_2$，需要单独考虑可能存在的直线斜率不存在的情况 \\
（2）$A, B, C$三点共线：$k_{AB} = k_{AC} = k_{BC}$ \\
（3）两直线垂直：直线斜率$k_1 \cdot k_2 = -1$，需要单独考虑可能存在的直线斜率不存在的情况；或者用向量的数量积等于0 \\
（4）两直线的倾斜角互余：直线斜率$k_1 \cdot k_2 = 1$，需要单独考虑可能存在的直线斜率不存在的情况 \\
（5）两直线的倾斜角互补、两直线关于$y$轴对称、两直线关于$x$轴对称：直线斜率$k_1 + k_2 = 0$ ，需要单独考虑可能存在的直线斜率不存在的情况\\
（6）两向量$\overrightarrow{a}$ 与 $\overrightarrow{b}$的夹角为锐角：$\overrightarrow{a}\cdot\overrightarrow{b}>0$，钝角反之\\
（7）等腰直角三角形 $ABC$：以 $\angle A$ 为直角为例，取斜边 $BC$ 的中点为 $D$，则$AB\perp AC,AD \perp BC$，或者$AB \perp AC, \vert AB \vert  = \vert AC \vert$\\
（8）直角三角形中的特殊角如$30^\circ$、$60^\circ$：以 $\angle A$ 为直角为例，若$\angle B = 30^\circ$，则$\tan B = \frac{\vert AC \vert }{\vert AB \vert}= \frac{\sqrt{3}}{3}$ \\


\begin{example} % 直线引入，直线相交求交点，三点共线
（2026朝阳一模）已知椭圆 $E:\frac{x^2}{a^2}+\frac{y^2}{b^2}=1\ (a>b>0)$ 的离心率为 $\frac{\sqrt{3}}{2}$，$A,B$ 分别为椭圆 $E$ 的上、下顶点，且 $B(0,-1)$.\\
（1）求椭圆 $E$ 的方程；\\
（2）过点 $D(-1,-1)$ 的斜率存在且不为 $1$ 的直线 $l$ 与椭圆 $E$ 交于不同的两点 $M,N$（均不与点 $A$ 重合），点 $P$ 与点 $M$ 关于原点 $O$ 对称，直线 $BP$ 与直线 $OD$ 交于点 $Q$. 求证：直线 $NQ$ 经过点 $A$.
\end{example}

\begin{example} % 动点引入，韦达定理求坐标，三点共线
（2026东城二模）已知椭圆 $W:\frac{x^2}{a^2}+\frac{y^2}{b^2}=1\ (a>b>0)$ 的右焦点为 $(\sqrt{3},0)$，且经过点 $\left(-1,\frac{\sqrt{3}}{2}\right)$。\\
（1）求椭圆 $W$ 的方程；\\
（2）已知点 $A(a,0),T(0,t)\ (t\neq\pm1)$。设直线 $x=a$ 分别与直线 $y=b,y=-b$ 交于点 $M,N$，直线 $MT$ 与直线 $y=-b$ 交于点 $P$，直线 $NT$ 与直线 $y=b$ 交于点 $Q$，直线 $AT$ 交椭圆 $W$ 于另一点 $R$。求证：$P,Q,R$ 三点共线。
\end{example}

\begin{example} % 直线，三点共线
（2026朝阳期末）已知椭圆 $E:\frac{x^2}{a^2}+\frac{y^2}{b^2}=1\ (a>b>0)$ 的离心率为 $\frac{\sqrt{2}}{2}$，椭圆 $E$ 上的点到两焦点的距离之和为 $2\sqrt{6}$.\\
（1）求椭圆 $E$ 的方程；\\
（2）设点 $M(m,0),N(n,0)\ (m\neq n)$，过点 $M$ 的直线 $l$ 与椭圆 $E$ 交于不同的两点 $A,B$，点 $A$ 关于 $x$ 轴的对称点记为 $P$（$P$ 与 $B$ 不重合）. 若 $mn=6$，试判断 $P,B,N$ 三点是否共线，并说明理由.
\end{example}

\begin{example} % 直线引入，倾斜角互余
（2025朝阳二模）已知椭圆 $E:\frac{x^2}{a^2}+\frac{y^2}{b^2}=1\ (a>b>0)$ 的焦距为 $2$，且过点 $(0,\sqrt{3})$。\\
（1）求椭圆 $E$ 的方程；\\
（2）设直线 $y=kx+2\ (k<0)$ 与椭圆 $E$ 交于不同的两点 $A,B$，直线 $y=x$ 与直线 $AB$ 交于点 $N$。若 $\angle AON=\angle BON$（$O$ 为坐标原点），求 $k$ 的值。
\end{example}

\begin{example} % 直线或者动点，距离公式，等腰直角三角形
（2026海淀期末）已知椭圆 $E:\frac{x^2}{a^2}+\frac{y^2}{b^2}=1\ (a>b>0)$ 的左顶点为 $A$，请从下面条件中任选两个作为已知.\\
条件①：椭圆 $E$ 的离心率为 $\frac{1}{2}$；\\
条件②：椭圆 $E$ 经过点 $(0,\sqrt{3})$；\\
条件③：点 $A$ 到椭圆 $E$ 的左焦点的距离为 $1$.\\
（1）求椭圆 $E$ 的方程；\\
（2）过点 $A$ 的直线 $l$ 与椭圆 $E$ 的另一个交点为 $P$，$l$ 与 $y$ 轴交于点 $M$. 过点 $P$ 作 $l$ 的垂线交直线 $x=2$ 于点 $N$. 若 $\triangle PMN$ 是等腰直角三角形，求 $l$ 的斜率.
\end{example}

\begin{example} % 弦长公式、动点在直线上，韦达定理求坐标
（2026西城期末）已知椭圆 $E:\frac{x^2}{2} + y^2 = 1$ 的上顶点为 $A$，动点 $B$ 为 $x$ 轴正半轴上的一点，过点 $A$ 作 $AB$ 的垂线交椭圆 $E$ 于另一点 $C$.\\
(1) 当直线 $AC$ 经过椭圆的左焦点时，求点 $B$ 的坐标；\\
(2) 是否存在点 $B$，分别使 $\angle ACB = 30^{\circ}$、$45^{\circ}$ 和 $60^{\circ}$？说明理由。 

\end{example}

\begin{example} 
（2025海淀二模）已知椭圆 $C:\frac{x^2}{6}+\frac{y^2}{2}=1$。设直线 $l:y=x+m$ 交椭圆 $C$ 于不同的两点 $A,B$，与 $y$ 轴交于点 $P$。\\
（1）当 $m=0$ 时，求 $\lvert AB\rvert$ 的值；\\
（2）若点 $Q$ 满足 $\lvert PQ\rvert=3$ 且 $\lvert QA\rvert=\lvert QB\rvert$，求 $\angle AQB$ 的大小。
\end{example}


\subsection{面积}
\noindent （1）三角形的面积公式：$S_{\triangle ABC}=\frac{1}{2}\cdot \text{底}\cdot \text{高}=\frac{1}{2}ab\sin C$。
优先考虑弦长公式配合第三点到直线的距离求面积；当已知两边及夹角时，直接使用 $\frac{1}{2}ab\sin C$。\\
（2）原点三角形的面积公式：
若$O$为坐标原点，$A(x_1,y_1),B(x_2,y_2)$，则
$S_{\triangle OAB}=\frac{1}{2}\left|x_1y_2-x_2y_1\right|$ \\
该公式可以扩展到任意三角形面积，如求$\triangle ABC$的面积，求出$\overrightarrow{AB},\overrightarrow{AC}$即可，然后套公式\\
（3）铅锤法 \\
给定一个$\triangle ABC$，过点$A$作平行于$x$轴的直线，与$BC$或者$BC$的延长线交于点$D_1$，则$S_{\triangle ABC} = \frac{1}{2} \cdot \vert x_A - x_{D_1} \vert \cdot \vert y_B - y_C \vert$；过点$A$作平行于$y$轴的直线，与$BC$或者$BC$的延长线交于点$D_2$，则$S_{\triangle ABC} = \frac{1}{2} \cdot \vert y_A - y_{D_2} \vert \cdot \vert x_B - x_C \vert$。同理可以用点$B,C$作相同的操作，一共可以获得6个公式。\\

\begin{center}
\begin{tikzpicture}[scale=0.9]
    \coordinate (A) at (1,3);
    \coordinate (B) at (0,0);
    \coordinate (C) at (4,4);
    \coordinate (D1) at (3,3);
    \coordinate (D2) at (1,1);

    \draw[thick] (A) -- (B) -- (C) -- cycle;
    \draw[blue, thick, dashed] (A) -- (D1);
    \draw[red, thick, dashed] (A) -- (D2);
    \fill (A) circle (1.5pt) node[above left] {$A$};
    \fill (B) circle (1.5pt) node[below left] {$B$};
    \fill (C) circle (1.5pt) node[above right] {$C$};
    \fill (D1) circle (1.5pt) node[above right] {$D_1$};
    \fill (D2) circle (1.5pt) node[below right] {$D_2$};
    \node[blue, above] at (2,3) {水平宽};
    \node[red, left] at (1,2) {竖直高};
\end{tikzpicture}
\end{center}
适用场景：作了水平宽或者竖直高之后，水平宽或者竖直高比较容易计算，如果还需要去额外计算交点的坐标的话，就不如直接用底乘高求面积了。\\
如设坐标原点为$O$，过$y$轴上一定点$(0, n)$的直线上两点$A(x_1, y_1), B(x_2, y_2)$，则$$S_{\triangle OAB} = \frac{1}{2} \cdot \vert n \vert \cdot \vert x_1 - x_2 \vert$$
同理，过$x$轴上一定点$(n, 0)$的直线上两点$A(x_1, y_1), B(x_2, y_2)$，则$$S_{\triangle OAB} = \frac{1}{2} \cdot \vert n \vert \cdot \vert y_1 - y_2 \vert$$ \\
（4）面积之比：同高三角形面积比等于底边比；同底三角形面积比等于高之比，即
\[
\frac{S_{\triangle ABC}}{S_{\triangle ABD}}=\frac{BC}{BD}\quad(B,C,D\text{ 三点共线}),\qquad
\frac{S_{\triangle ABC}}{S_{\triangle BCD}}=\frac{d(A,BC)}{d(D,BC)}.
\]
若两个三角形有相等的角或者互补的角，则可以根据$S = \frac{1}{2}ab \sin C$将面积比转化为该角两边的长度乘积之比。


\begin{example} % 直线引入，三角形面积
（2025东城二模）已知椭圆 $E:\frac{x^2}{a^2}+\frac{y^2}{b^2}=1\ (a>b>0)$ 的一个顶点为 $A(0,1)$，且过点 $\left(1,\frac{\sqrt{3}}{2}\right)$。\\
（1）求椭圆 $E$ 的方程及焦距；\\
（2）过点 $(0,2)$ 的直线与椭圆 $E$ 交于不同的两点 $B,C$。直线 $AB,AC$ 的斜率分别记为 $k_1,k_2$，当 $k_1+k_2=\frac{2}{3}$ 时，求 $\triangle ABC$ 的面积。
\end{example}

\begin{example}
（2026北京高考）已知椭圆 $E:\frac{x^2}{a^2}+\frac{y^2}{b^2}=1\ (a>b>0)$ 的一个顶点为 $(2,0)$，离心率为 $\frac{1}{2}$。\\
（1）求椭圆 $E$ 的方程；\\
（2）过点 $A(1,1)$，斜率为 $k\ (k\neq\pm1)$ 的直线交椭圆 $E$ 于 $B,C$ 两点，点 $B$ 关于 $y=x$ 的对称点为 $D$，直线 $DC$ 与 $y=x$ 交于点 $Q$。若
$\left|S_{\triangle ABQ}-S_{\triangle ACQ}\right|=\frac{5}{8}$，求 $k$ 的值。
\end{example}

\begin{example}
（2025北京高考）已知椭圆 $E:\frac{x^2}{a^2}+\frac{y^2}{b^2}=1\ (a>b>0)$ 的离心率为 $\frac{\sqrt{2}}{2}$，椭圆 $E$ 上的点到两焦点的距离之和为 $4$。\\
（1）求椭圆 $E$ 的方程；\\
（2）设 $O$ 为坐标原点，点 $M(x_0,y_0)\ (x_0\neq0)$ 在椭圆 $E$ 上，直线 $x_0x+2y_0y-4=0$ 与直线 $y=2$、$y=-2$ 分别交于点 $A,B$。设 $\triangle OAM$ 与 $\triangle OBM$ 的面积分别为 $S_1,S_2$，比较 $\frac{S_1}{S_2}$ 与 $\frac{\lvert OA\rvert}{\lvert OB\rvert}$ 的大小。
\end{example}


\begin{example}
（2026海淀二模）椭圆 $E:\frac{x^2}{a^2}+\frac{y^2}{b^2}=1\ (a>b>0)$ 的下顶点为 $A(0,-1)$，左焦点到点 $A$ 的距离为 $\sqrt{3}$。\\
（1）求椭圆 $E$ 的方程及离心率；\\
（2）经过点 $A$ 的直线与椭圆的另一个交点为 $D$（位于 $x$ 轴上方），与 $x$ 轴的交点为 $M$。点 $D$ 关于 $x$ 轴的对称点为 $E$，过点 $D$ 作与 $x$ 轴平行的直线，交直线 $AE$ 于点 $N$。若
$S_{\triangle DEN}=3S_{\triangle AEM}$，求直线 $AD$ 的斜率。
\end{example}

\begin{example} % 直线引入，铅锤法求面积
（2023西城二模）已知椭圆 $E:\frac{x^2}{a^2}+\frac{y^2}{b^2}=1\ (a>b>0)$ 的短轴长为 $2\sqrt{2}$，一个焦点为 $F_1(-2,0)$。\\
（1）求椭圆 $E$ 的方程和离心率；\\
（2）设直线 $l:x-my-2=0$ 与椭圆 $E$ 交于两点 $A,B$，点 $M$ 在线段 $AB$ 上，点 $F_1$ 关于点 $M$ 的对称点为 $C$。当四边形 $AF_1BC$ 的面积最大时，求 $m$ 的值。
\end{example}





\subsection{圆}
\noindent（1）点$P$在以$AB$为直径的圆上、圆内、圆外：$ <\overrightarrow{PA}, \overrightarrow{PB}>$等于、大于、小于$90^\circ$ \\
（2）圆的切线：点到直线的距离等于半径；圆心与切点的所成的直线垂直于切线




\begin{example} % 动点、以线段为直径的圆过某点
（2025东城一模）已知椭圆 $E:\frac{x^{2}}{a^{2}}+\frac{y^{2}}{b^{2}}=1\ (a>b>0)$ 过点 $(0,1)$，离心率为 $\frac{\sqrt{6}}{3}$，$E$ 上的点 $A(m,n)\ (n\neq0)$ 关于 $x$ 轴的对称点为 $B$. 设 $O$ 为原点，$\overrightarrow{OH}=\lambda\overrightarrow{OA}\ (0<\lambda<1)$，过点 $H$ 与 $x$ 轴平行的直线交 $E$ 于点 $P,Q$.\\
(1) 求椭圆 $E$ 的方程；\\
(2) 若点 $B$ 在以 $PQ$ 为直径的圆上，求 $\lambda$ 的值.
\end{example}


\begin{example} % 动点引入，椭圆与圆的关系，基本不等式求最值
（2022西城二模）已知椭圆 $C:\frac{x^2}{a^2}+\frac{y^2}{b^2}=1\ (a>b>0)$ 的左顶点为 $A(-2,0)$，圆 $O:x^2+y^2=1$ 经过椭圆 $C$ 的上、下顶点。\\
（1）求椭圆 $C$ 的方程和焦距；\\
（2）已知 $P,Q$ 分别是椭圆 $C$ 和圆 $O$ 上的动点（$P,Q$ 不在坐标轴上），且直线 $PQ$ 与 $x$ 轴平行，线段 $AP$ 的垂直平分线与 $y$ 轴交于点 $M$，圆 $O$ 在点 $Q$ 处的切线与 $y$ 轴交于点 $N$。求线段 $MN$ 长度的最小值。
\end{example}



\subsection{平行四边形}
\noindent（1）证明一个四边形$ABCD$为平行四边形：证明$\overrightarrow{AB} = \overrightarrow{DC} \text{ 或 } \overrightarrow{AB} + \overrightarrow{AD} = \overrightarrow{AC}$\\
（2）证明一个四边形为矩形、菱形、正方形：在证明平行四边形的基础分别加上邻边垂直、对角线垂直、邻边垂直且对角线垂直\\
（3）平行四边形$ABCD$的面积：$S_{ABCD} = 2 S_{\triangle ABC}$


\begin{example} % 直线、平行四边形翻译
（2026东城一模）已知椭圆 $C:\frac{x^2}{a^2}+\frac{y^2}{b^2}=1\ (a>b>0)$ 经过点 $(0,1)$，离心率为 $\frac{\sqrt{6}}{3}$。\\
（1）求椭圆 $C$ 的方程；\\
（2）过点 $D(1,0)$ 作斜率为 $k$ 的直线 $l$，与椭圆 $C$ 交于 $A,B$ 两点。点 $E(3,n)$ 不在直线 $l$ 上，直线 $AE,BE$ 分别与直线 $x=2$ 交于点 $M,N$。求证：四边形 $DMEN$ 为平行四边形。
\end{example}



\begin{example} % 直线引入，菱形的翻译
（2024海淀二模）已知椭圆 $E$ 的焦点在 $x$ 轴上，中心在坐标原点。以 $E$ 的一个顶点和两个焦点为顶点的三角形是等边三角形，且其周长为 $6\sqrt{2}$。\\
（1）求椭圆 $E$ 的方程；\\
（2）设过点 $M(2,0)$ 的直线 $l$（不与坐标轴垂直）与椭圆 $E$ 交于不同的两点 $A,C$，与直线 $x=16$ 交于点 $P$。点 $B$ 在 $y$ 轴上，$D$ 为坐标平面内的一点，四边形 $ABCD$ 是菱形。求证：直线 $PD$ 过定点。
\end{example}

\begin{example} % 直线、正方形翻译
（2025西城期末）已知椭圆 $E:\frac{x^2}{a^2}+\frac{y^2}{b^2}=1\ (a>b>0)$ 的左、右顶点分别为 $A_1,A_2$，离心率为 $\frac{\sqrt{2}}{2}$，点 $T(0,1)$，$\triangle TA_1A_2$ 的面积为 $2$。\\
（1）求椭圆 $E$ 的方程；\\
（2）过点 $T$ 且斜率为 $k$ 的直线交椭圆 $E$ 于点 $C,D$，线段 $CD$ 的垂直平分线交 $y$ 轴于点 $Q$，点 $Q$ 关于直线 $CD$ 的对称点为 $P$。若四边形 $PCQD$ 为正方形，求 $k$ 的值。
\end{example}




\subsection{椭圆的切线}
\noindent 当题目中出现椭圆的切线时，通常有两种处理方式：\\
（1）设直线$y = kx + m$，将其与椭圆方程联立，令$\Delta = 0$，得到$k, m$之间的关系，然后进一步计算得到切点坐标\\
（2）设切线为$P(x_0, y_0)$，写出直线的点斜式方程，联立椭圆方程，令$\Delta = 0$，得到切线的斜率
\begin{example} % 直线动点，椭圆切线
（2024丰台期末）已知椭圆 $E:\frac{x^2}{4}+\frac{y^2}{3}=1$.\\
（1）求椭圆 $E$ 的离心率和焦点坐标；\\
（2）设直线 $l_1:y=kx+m$ 与椭圆 $E$ 相切于第一象限内的点 $P$，不过原点 $O$ 且平行于 $l_1$ 的直线 $l_2$ 与椭圆 $E$ 交于不同的两点 $A,B$，点 $A$ 关于原点 $O$ 的对称点为 $C$. 记直线 $OP$ 的斜率为 $k_1$，直线 $BC$ 的斜率为 $k_2$，求 $\frac{k_1}{k_2}$ 的值.
\end{example}

\begin{example} % 椭圆切线，距离公式
（2021东城期末）已知椭圆 $C:\frac{x^2}{a^2}+\frac{y^2}{b^2}=1\ (a>b>0)$ 过点 $A(-2,0),B(2,0)$，且离心率为 $\frac{1}{2}$.\\
（1）求椭圆 $C$ 的方程；\\
（2）设直线 $l$ 与椭圆 $C$ 有且仅有一个公共点 $E$，且与 $x$ 轴交于点 $G$（$E,G$ 不重合），$ET\perp x$ 轴，垂足为 $T$. 求证：$\frac{\lvert TA\rvert}{\lvert TB\rvert}=\frac{\lvert GA\rvert}{\lvert GB\rvert}$.
\end{example}




\newpage


